\chapter{Stability of the Explicit Field Update}
\label{app:stability}

% Source: ansys_usermat/apdl/nem_stability.py, results/2026-10-wp3_fix/nem_stability.json
% (9 Oct 2026). Deck values of the element: NEIGHBOR_CNT = 30, BETA_STAR = 0.2 mm,
% WMAT_THRESHOLD = 1e-8 (OLIVER_MODEL_NOTES.md).

The element advances its field explicitly,
$\phi_{n+1}=\phi_n+\Delta t\,(\beta L\phi_n+\ldots)$, where $L$ is the
Neighbored Element Method approximation of the Laplacian
\cite{Blaszczyk2022TaylorWLS,Rudolf2025NEM}: at every integration point a
second-order Taylor polynomial is fitted by weighted least squares through
the values at the $N$ nearest integration points, with the weight
$w=\exp[-\tfrac12(4d/\beta^*)^2]$ of the distance $d$, and the Laplacian is
the sum of the three unmixed second derivatives of the fit. In the runs of
Chapter~\ref{ch:ansys} the element uses $N=30$ and $\beta^*=0.2$~mm. For
explicit Euler the update is stable if $\beta\Delta t\,\rho\le2$, with $\rho$
the spectral radius of $-L$; written with $\lambda=\beta\Delta t/h^2$ this is
$\lambda\le\lambda_{\max}=2/(\rho h^2)$. For the three-point stencil in one
dimension $\rho=4/h^2$ and $\lambda_{\max}=1/2$, which is the condition
$h^2/(2\Delta t)\ge\beta$ of Rudolf et al.~\cite{Rudolf2025NEM}; for the
five-point stencil in two dimensions $\lambda_{\max}=1/4$ and for the
seven-point stencil in three dimensions $1/6$.

\paragraph{Lattice of the integration points.} On a uniform hexahedral mesh
of size $h$ the $2\times2\times2$ integration points of SOLID185 lie, per
axis, at $h(\tfrac12\pm\tfrac{1}{2\sqrt3})$ inside each element, so their
spacing alternates between $h/\sqrt3\approx0.58h$ inside an element and
$h(1-1/\sqrt3)\approx0.42h$ across the element boundary. The eight points of
an element are equivalent under the reflections of the cube, so the stencil
of one of them gives the other seven by mirroring, and the operator $L$ is
periodic with the period $h$ of the mesh with eight points per cell.

\paragraph{Spectral radius.} I computed the stencil of the WLS fit on this
lattice with the deck values and formed the $8\times8$ Bloch symbol
$\hat L(\mathbf k)$ of $L$ over the Brillouin zone of the cell lattice; $\rho$
is the largest modulus of its eigenvalues over $\mathbf k$. Two points
follow from the result (Table~\ref{tab:app-stability}).

First, the limit of the NEM stencil differs from $1/6$ even without the
weighting: with $\beta^*\to\infty$ it is $0.18$ to $0.19$, because the
lattice is not uniform and the stencil reaches beyond the nearest neighbours.

Second, the weighting makes the limit depend on the mesh. Because
$\beta^*$ is a length in millimetres, the weight of a neighbour at distance
$d$ depends on $d/\beta^*$, and on a coarse mesh the fit is dominated by the
few nearest points, which gives a stencil with large coefficients and a small
$\lambda_{\max}$: about $0.06$ on $8^3$ elements, $0.14$ to $0.15$ on $16^3$
and $0.16$ to $0.17$ on $24^3$ (the two values are for $N=30$ and for the
symmetric set of $32$ neighbours, between which the distance ordering of the
element does not decide). The limit approaches the unweighted value as the
mesh is refined. A direct time stepping of the stencil on a periodic block of
elements grows by the factor $|1-2\lambda/\lambda_{\max}|$ per step, as the
analysis predicts.

\begin{figure}[htbp]
  \centering
  \includegraphics[width=0.75\linewidth]{\ansysfigdir fig_nem_stability.png}
  \caption[Stability limit of the NEM stencil against the mesh]{Stability
    limit of the explicit field update from the Bloch analysis, against the
    number of elements per edge, for $30$, $32$ and $56$ neighbours with
    $\beta^*=0.2$~mm; bars: the runs of the element, from the largest smooth
    to the smallest oscillating $\lambda$ (Figure~\ref{fig:ch5-stability}).}
  \label{fig:app-stability}
\end{figure}

\paragraph{Comparison with the runs.} The runs of the element bracket the
limit at $0.128<\lambda_{\max}<0.160$ on $16^3$ and at
$0.158<\lambda_{\max}<0.187$ on $24^3$ (Figure~\ref{fig:ch5-stability}), in
agreement with the computed values. On $8^3$ no run of this work exceeds
$\lambda=0.04$, so the predicted limit of about $0.06$ there is not yet
tested. The unstable mode on the finer meshes is the alternation of sign
between neighbouring points, which is what the oscillating runs show in
$\alpha-1$.

\begin{table}[htbp]
  \centering
  \caption[Stability limit of the NEM stencil]{Stability limit
    $\lambda_{\max}=\beta\Delta t/h^2$ of the explicit update with the NEM
    stencil on the integration-point lattice, from the Bloch analysis;
    $\beta^*=0.2$~mm as in the element, and the unweighted limit. Last row:
    the largest smooth and the smallest oscillating $\lambda$ of the runs.}
  \label{tab:app-stability}
  \small
  \begin{tabular}{lcccc}
    \toprule
    & $8^3$ & $16^3$ & $24^3$ & $\beta^*\to\infty$ \\
    \midrule
    $N=30$ neighbours & $0.062$ & $0.139$ & $0.162$ & $0.185$ \\
    $N=32$ neighbours & $0.062$ & $0.149$ & $0.173$ & $0.180$ \\
    runs: smooth / oscillating & $\le0.04$ / -- & $0.128$ / $0.160$ & $0.158$ / $0.187$ & \\
    \bottomrule
  \end{tabular}
\end{table}
